\chapter{Sequences of Real Numbers}
\label{ch:sequences}

Threshold indices will usually be written explicitly in this chapter. A
direct proof must exhibit which choices work together; later applications
can invoke a proved theorem that already contains this bookkeeping.
Discovering a proof and writing it are different tasks:
\[
\boxed{\text{Discovery: reason backward from the desired error bound.}}
\]
\[
\boxed{\text{Proof: choose the threshold and verify forward.}}
\]
For $\frac1n\to0$, the target $\frac1n<\eps$ suggests
$n>\frac1\eps$.
The proof starts instead by choosing a natural number $N>1/\eps$ and
then checks $n\geq N\Rightarrow \frac1n\leq\frac1N<\eps$.
For $|a_n-a|+|b_n-b|<\eps$, an \emph{error budget} allocates
$\eps/2$ to each summand. Different allocations are possible; each
must be chosen before verifying the final estimate.

\section{Convergence}

We now study systematically the convergence used in Chapter~\ref{ch:reals},
beginning by restating its precise error criterion.
The error must eventually stay below each prescribed tolerance; merely
returning close to a proposed limit infinitely often is not enough.


In this chapter the real numbers are the ordered field constructed in Chapter~\ref{ch:reals}. A sequence in $\R$ is a function $a:\N\to\R$; its value at $n$ is written $a_n$.

\begin{definition}[Limit of a sequence]
\label{def:sequence-limit}
Let $(a_n)$ be a real sequence and let $L\in\R$. We say that $(a_n)$ \textbf{converges to} $L$, and write
\[
\lim_{n\to\infty}a_n=L,
\]
or equivalently $a_n\to L$, if for every $\eps>0$ there exists $N\in\N$ such that, for every $n\geq N$, one has $\abs{a_n-L}<\eps$. A sequence that converges to some real number is \textbf{convergent}; otherwise it is \textbf{divergent}.
\end{definition}

The number $N$ may depend on $\eps$, but it must work for every $n\geq N$. This order of quantifiers is essential.

\subsection*{Revisiting decimals through sequence limits}
Section~\ref{sec:decimal-expansions} already interpreted decimals through
finite truncations and shrinking intervals, including the terminating
versus repeating-$9$ ambiguity. We now revisit those conclusions in the
language of \cref{def:sequence-limit}. The same three questions become
convergence of prescribed truncations, construction of digits for a given
real number, and uniqueness apart from endpoint ambiguity. Finite tail
estimates make the convergence explicit; no infinite-series theory is needed.
\begin{example}[Constructing digits by successive intervals]
For $x=1/7$, integer division successively gives
$10=1\cdot7+3$, $30=4\cdot7+2$, $20=2\cdot7+6$,
$60=8\cdot7+4$, $40=5\cdot7+5$, and $50=7\cdot7+1$.
Thus the digits repeat $142857$. The first three truncations are
$0.1,0.14,0.142$, and the $n$th truncation $d_n$ satisfies
$d_n\leq1/7<d_n+10^{-n}$. The shrinking error, rather than the
repeating pattern, gives $d_n\to1/7$ by the definition of convergence,
recovering Chapter~\ref{ch:reals}'s interval interpretation.
\end{example}
\begin{proposition}[Decimal representations]
\label{prop:decimal-representations}
For digits $d_j\in\{0,\ldots,9\}$, the truncations
$s_n=\sum_{j=1}^n d_j10^{-j}$ converge to a real number in $[0,1]$,
denoted $0.d_1d_2\cdots$.
Every $x\in[0,1)$ has such an expansion which is not eventually $9$.
Two different digit strings represent the same number only when, at their
first differing place, their digits differ by one and the smaller digit is
followed entirely by $9$'s, the larger entirely by $0$'s.
\end{proposition}
\begin{proof}
For $m>n$, a finite geometric sum gives
\[
0\leq s_m-s_n\leq\sum_{j=n+1}^m9\,10^{-j}
=10^{-n}-10^{-m}<10^{-n}.
\]
Given $\eps>0$, choose $N$ with $10^{-N}<\eps$ (use $10^N\geq N+1$).
For $m>n\geq N$,
$|s_m-s_n|<10^{-n}\leq10^{-N}<\eps$; swapping $m,n$ covers the other
order, and equal indices give zero. The truncations are Cauchy and converge by \cref{thm:real-cauchy-complete}.
Their bounds $0\leq s_n\leq1$ also bound their limit: a limit outside that
interval would have a positive distance from every truncation.

To construct digits for $x<1$, set $r_0=x$ and recursively choose the
integer $d_n$, using \cref{lem:integer-part}, with $d_n\leq10r_{n-1}<d_n+1$, putting
$r_n=10r_{n-1}-d_n$. Then $0\leq r_n<1$, $d_n\in\{0,\ldots,9\}$,
and $x=s_n+10^{-n}r_n$. For justification set $s_0=0$; the identity
holds at $n=0$. Substituting $r_{n-1}=(d_n+r_n)/10$ into the
identity at $n-1$ proves it at $n$. For $n\geq N$ as above,
$|x-s_n|=10^{-n}r_n<10^{-n}\leq10^{-N}<\eps$, so $s_n\to x$. This expansion cannot have
only $9$'s after position $n$: that tail would force $r_n=1$.

For uniqueness, suppose two strings first differ at place $k$, with
$a_k<b_k$. For every $m>k$, the maximum compensation from the next
$m-k$ digits is $10^{-k}-10^{-m}$ by the finite identity above.
The limiting compensation is therefore at most $10^{-k}$.
Equality of the represented numbers
therefore forces $b_k-a_k=1$ and the maximal compensating tails. If even
one later digit fails to be maximal, the compensation loses a fixed positive
amount. Thus the tails must be respectively all $9$'s and all $0$'s.
\end{proof}

The endpoint example from \cref{sec:decimal-expansions} now reads
\[
0.\underbrace{99\cdots9}_{n\text{ digits}}=1-10^{-n},
\qquad 0.999\cdots=1.
\]
Likewise $0.24999\cdots=0.25$. A decimal expression represents a number;
it is not the number itself. For a canonical representation use the one
that is not eventually $9$ (and the usual integer part for numbers outside
$[0,1)$).

\begin{theorem}[Cantor's diagonal argument, now formal]
\label{thm:cantor-formal}
The interval $(0,1)$, and hence $\R$, is uncountable.
\end{theorem}
\begin{proof}
If $(x_n)$ listed $(0,1)$, write each entry in its canonical decimal
expansion. Choose the $n$th digit $c_n$ to be $2$ if the $n$th digit of
$x_n$ is $1$, and $1$ otherwise. The proposition gives a real
number $c=0.c_1c_2\cdots$ in $(0,1)$. Its digits are only $1$ and $2$,
so no terminating/repeating-$9$ ambiguity can identify it with an $x_n$
whose $n$th digit differs. This contradicts the claimed listing.
\end{proof}
This completes the informal diagonal-argument preview in
Chapter~\ref{ch:foundations}, using the decimal interpretation established
in Chapter~\ref{ch:reals} and made explicit through sequence limits here.

\begin{lemma}
\label{lem:positive-small}
If $x\in\R$ satisfies $\abs{x}<\eps$ for every $\eps>0$, then $x=0$.
\end{lemma}
\begin{proof}
If $x\ne0$, then $\abs{x}>0$. Taking $\eps=\abs{x}/2$ contradicts the hypothesis.
\end{proof}

\begin{theorem}[Uniqueness of limits]
\label{thm:unique-sequence-limit}
A convergent real sequence has exactly one limit.
\end{theorem}
\begin{proof}
Suppose $a_n\to L$ and $a_n\to M$. Given $\eps>0$, choose $N_1,N_2$ such that $\abs{a_n-L}<\eps/2$ for $n\geq N_1$ and $\abs{a_n-M}<\eps/2$ for $n\geq N_2$. For $n\geq\max\{N_1,N_2\}$, the triangle inequality gives
\[
\abs{L-M}\leq\abs{L-a_n}+\abs{a_n-M}<\eps.
\]
By \cref{lem:positive-small}, $L-M=0$.
\end{proof}

\begin{example}
The sequence $a_n=1/n$ converges to $0$. For a given $\eps>0$, \cref{thm:archimedean} provides $N>1/\eps$, and then $n\geq N$ implies
\[
0<\frac1n\leq\frac1N<\eps.
\]
\end{example}

\begin{definition}[Divergence to infinity]
\label{def:sequence-infinite-limit}
We write
\[
a_n\longrightarrow+\infty
\]
if, for every $M\in\R$, there is $N\in\N$ such that
\[
n\geq N\quad\Longrightarrow\quad a_n>M.
\]
We write $a_n\to-\infty$ if, for every $M\in\R$, there is $N$ such that
\[
n\geq N\quad\Longrightarrow\quad a_n<M.
\]
These statements describe two particular forms of divergence.  Neither
$+\infty$ nor $-\infty$ is a real number, and neither statement says that
$(a_n)$ converges in $\R$.
\end{definition}

\begin{example}[Examples and nonexamples]
The sequence $n^2$ tends to $+\infty$: given $M\in\R$, choose
$N>\max\{M,1\}$.  The sequence $-n$ tends to $-\infty$.  By contrast,
$(-1)^n n$ is unbounded in both directions but tends to neither infinity;
for every proposed stage, later terms occur with each sign.  Thus
``unbounded'' is much weaker than ``eventually larger than every real
bound.''
\end{example}

\begin{remark}[A useful proof pattern]
To prove $a_n\to+\infty$, start with an arbitrary target height $M$ and
solve the inequality $a_n>M$ for a sufficient condition on $n$.  This is
the infinite-limit analogue of solving $\abs{a_n-L}<\eps$ in an ordinary
$\eps$--$N$ proof.
\end{remark}

\section{Algebra and Order of Limits}

Once limits have been defined, the next question is whether familiar algebra
survives passage to the limit.  The following rules justify the calculations
that will later occur inside functions, series, derivatives, and integrals.

\begin{theorem}[Algebra of limits]
\label{thm:limit-algebra}
Let $(a_n)$ and $(b_n)$ be real sequences with $a_n\to a$ and $b_n\to b$, and let $c\in\R$. Then
\[
\lim_{n\to\infty}(a_n+b_n)=a+b,
\qquad
\lim_{n\to\infty}ca_n=ca,
\qquad
\lim_{n\to\infty}a_nb_n=ab.
\]
If $b\ne0$, then
\[
\lim_{n\to\infty}\frac{a_n}{b_n}=\frac ab
\]
whenever $b_n\ne0$ for all sufficiently large $n$.
\end{theorem}
\begin{proof}
Fix $\eps>0$. For sums choose $N_1,N_2$ so that
$|a_n-a|<\eps/2$ for $n\geq N_1$ and $|b_n-b|<\eps/2$ for
$n\geq N_2$. Put $N=\max\{N_1,N_2\}$. For every $n\geq N$,
\[
|(a_n+b_n)-(a+b)|\leq|a_n-a|+|b_n-b|<\eps.
\]
If $c=0$, scalar multiplication has zero error. Otherwise choose $N_c$
with $|a_n-a|<\eps/|c|$ for $n\geq N_c$; then
$|ca_n-ca|=|c||a_n-a|<\eps$ for every such $n$.

For products first choose $N_0$ with $|a_n-a|<1$ for $n\geq N_0$.
Set $B=|a|+1$; thus $|a_n|<B$ on that tail. Choose $N_1$ for
$|b_n-b|<\eps/(2B)$ and $N_2$ for $|a_n-a|<\eps/(2(|b|+1))$.
For every $n\geq N=\max\{N_0,N_1,N_2\}$,
\[
|a_nb_n-ab|\leq |a_n||b_n-b|+|b||a_n-a|<\eps.
\]
The bound $B$ was fixed before choosing the error thresholds.

For quotients choose $N_0$ with $|b_n-b|<|b|/2$ for $n\geq N_0$.
Then $|b_n|\geq|b|-|b_n-b|>|b|/2>0$, so the quotient is defined.
Choose $N_1$ for $|a_n-a|<\eps|b|/4$ and $N_2$ for
$|b_n-b|<\eps|b|^2/(4(|a|+1))$. For every
$n\geq N=\max\{N_0,N_1,N_2\}$,
\[
\left|\frac{a_n}{b_n}-\frac ab\right|
\leq\frac2{|b|}|a_n-a|+\frac{2|a|}{|b|^2}|b_n-b|<\eps.
\]
Thus the denominator hypothesis on the tail follows from $b\ne0$ itself.
\end{proof}

\begin{theorem}[Order preservation]
\label{thm:limit-order}
Let $a_n\to a$ and $b_n\to b$. If $a_n\leq b_n$ for every sufficiently large $n$, then $a\leq b$.
\end{theorem}
\begin{proof}
Let $N_0$ be a stage from which $a_n\leq b_n$. If $a>b$, put
$\eps=(a-b)/3$. Choose $N_1,N_2$ for the respective errors
$|a_n-a|<\eps$ and $|b_n-b|<\eps$. For every
$n\geq\max\{N_0,N_1,N_2\}$,
$a_n>a-\eps=b+2\eps>b_n$, a contradiction.
\end{proof}

\begin{corollary}[Squeeze theorem]
\label{cor:squeeze}
If $a_n\leq b_n\leq c_n$ eventually and $a_n\to L$, $c_n\to L$, then $b_n\to L$.
\end{corollary}
\begin{proof}
Let $N_0$ be a stage from which the ordering holds. Given $\eps>0$,
choose $N_1,N_2$ for $|a_n-L|<\eps$ and $|c_n-L|<\eps$,
respectively. Put $N=\max\{N_0,N_1,N_2\}$. For every $n\geq N$,
$L-\eps<a_n\leq b_n\leq c_n<L+\eps$, so $|b_n-L|<\eps$.
\end{proof}


\begin{example}[Integer powers and geometric sequences]
A geometric sequence with initial term $a$ and common ratio $r$ is
$a,ar,ar^2,\ldots$. We first study powers of the ratio.
For fixed $n\in\N$, the estimate $0<\frac1{k^n}\leq\frac1k$ holds
for $k\geq1$. Given $\eps>0$, choose $N>\max\{1,1/\eps\}$; every
$k\geq N$ then has $\frac1{k^n}<\eps$. Here $n$ is fixed and $k$ varies.
For $r>1$, Bernoulli's inequality supplies the growing lower bound
\[
r^n=(1+(r-1))^n\geq1+n(r-1).
\]
Given $M\in\R$, choose $N>\max\{1,(M-1)/(r-1)\}$. Then $n\geq N$
implies $r^n>M$, proving divergence to $+\infty$.
If $0<|r|<1$, set $h=\frac1{|r|}-1>0$. The same inequality gives
\[
|r^n|\leq\frac1{1+nh}.
\]
Choosing $N>\max\{1,1/(h\eps)\}$ now makes $|r^n|<\eps$ for every
$n\geq N$. For $r=0$, all positive powers are zero.
\end{example}

\begin{proposition}[Dominant powers in polynomial substitutions]
\label{prop:polynomial-sequence}
Let $f(x)=a_nx^n+\cdots+a_0$ and $g(x)=b_mx^m+\cdots+b_0$
be nonzero polynomial expressions, with $a_n\ne0$, $b_m\ne0$,
$\deg f=n$ and $\deg g=m$ (degrees may be zero).
Here $x$ is an indeterminate; substituting natural numbers gives the
sequence $f(k)/g(k)$ wherever its denominator is nonzero.
From some stage onward this denominator is nonzero, and
\[
\frac{f(k)}{g(k)}\longrightarrow
\begin{cases}
0,&n<m,\\
\dfrac{a_n}{b_m},&n=m,\\
+\infty,&n>m,\ \dfrac{a_n}{b_m}>0,\\
-\infty,&n>m,\ \dfrac{a_n}{b_m}<0.
\end{cases}
\]
\end{proposition}
\begin{proof}
Divide by the leading powers, not by a guessed limit:
\[
F_k=\frac{f(k)}{k^n}=a_n+\frac{a_{n-1}}k+\cdots+\frac{a_0}{k^n}\to a_n,
\qquad G_k=\frac{g(k)}{k^m}\to b_m.
\]
For degree zero the corresponding expression is just its constant term.
These limits follow from the integer-power example and limit arithmetic.
Choose $N_0$ such that $|G_k-b_m|<|b_m|/2$ for every $k\geq N_0$.
Then $|G_k|>|b_m|/2$, so $g(k)\ne0$ there, and
$H_k=F_k/G_k\to c=a_n/b_m$ by the quotient theorem.
Now $f(k)/g(k)=k^{n-m}H_k$. For $n<m$, write this as
$H_k/k^{m-n}$ and use limit arithmetic; for $n=m$ it is $H_k$.
For $n>m$ and $c>0$, choose $N_1$ such that $|H_k-c|<c/2$ for
$k\geq N_1$. Given $M\in\R$, choose a natural number
\[
N\geq\max\{N_0,N_1\},
\qquad
N>\max\left\{1,\frac{2\max\{M,0\}}c\right\}.
\]
Every $k\geq N$ satisfies $k^{n-m}H_k>(c/2)k>M$.
For $c<0$, apply this positive argument to $-f(k)/g(k)$, using
target $-M$ to obtain $f(k)/g(k)<M$.
\end{proof}
For example, $(2k+1)/(k^2+1)\to0$,
$(3k^2-k)/(2k^2+5)\to3/2$, and
$(-k^3+2)/(k+1)\to-\infty$. The leading coefficients decide the
sign in the last case; the lower terms cannot reverse it on the controlled tail.

\begin{example}[Algebraic simplification and rationalization]
The candidate limit in
\[
a_n=\frac{3n^2-n}{2n^2+5}
\]
becomes visible after division by $n^2$; the limit laws then give
$a_n\to3/2$.  For
\[
b_n=\sqrt{n^2+n}-n,
\]
direct substitution hides the answer.  Multiplying by the conjugate gives
\[
b_n=\frac{n}{\sqrt{n^2+n}+n}
=\frac{1}{\sqrt{1+\frac1n}+1}\longrightarrow\frac12.
\]
Here square roots exist by \cref{prop:positive-square-roots}; their needed
limit follows directly from
\[
0\leq\sqrt{1+\frac1n}-1
=\frac{\frac1n}{\sqrt{1+\frac1n}+1}
\leq\frac1n.
\]
The algebra does not itself prove a limit; it rewrites the expression into
terms to which already proved limit laws apply.
\end{example}

\begin{remark}[Strategy for sequence limits]
Separate discovery from proof.  First simplify, compute a few terms, or
compare dominant terms to guess a candidate.  Then prove it using the
definition, the limit laws, a squeeze estimate, or monotonicity and
boundedness.  To prove that no finite limit exists, look for incompatible
subsequences, failure of the Cauchy condition, or divergence to
$\pm\infty$.  A failed attempt to find a limit is not a proof of
nonexistence.
\end{remark}


\begin{example}[Translating repeated decay]
A process removes half the remaining material at each step, starting
with $8$ grams. The varying quantity is the amount $a_n$ after $n$
steps; the assumption is that the same fraction is removed each time
and none is added. Thus
\[
a_0=8,
\qquad
a_{n+1}=\frac{a_n}{2}.
\]
Induction gives the model
\[
a_n=8\left(\frac12\right)^n.
\]
The geometric-sequence
result proves $a_n\to0$: any positive tolerance in grams is met
and retained after some step. For error below $1/10$ gram, $N=7$
works, since $8/2^7=1/16<1/10$ and the amounts decrease.
Convergence does not mean the amount becomes exactly zero at a finite step.
\end{example}

\paragraph{Disproving a proposed limit.}
Negating $a_n\to L$ gives precisely
\[
\exists\eps_0>0\ \forall N\in\N\ \exists n\geq N:
\quad |a_n-L|\geq\eps_0.
\]
For $a_n=1+1/n$, the candidate $L=0$ fails: take $\eps_0=1$ and
$n=N$. This does not prove divergence; in fact $a_n\to1$.
To prove divergence one must exclude \emph{every} real $L$.
For $(-1)^n$, given $L\geq0$ choose any odd $n\geq N$;
given $L<0$ choose any even $n\geq N$. In either case
$|(-1)^n-L|\geq1$. The next subsection packages this obstruction
into the particularly useful two-subsequence test.

\section{Subsequences and Elementary Divergence}

A subsequence samples an infinite process along a thinner infinite set of
stages.  It is especially useful for detecting divergence: incompatible
long-term behaviors cannot belong to one convergent sequence.

\begin{definition}[Subsequence]
Let $(a_n)$ be a sequence. A \textbf{subsequence} is a sequence $(a_{n_k})_{k\in\N}$, where $n_1<n_2<\cdots$ are natural numbers.
\end{definition}

\begin{proposition}[Subsequences preserve limits]
\label{prop:subsequence-limit}
If $a_n\to L$, then every subsequence of $(a_n)$ converges to $L$.
\end{proposition}
\begin{proof}
For $n_1<n_2<\cdots$, induction gives $n_k\geq k$. Given $\eps>0$, choose $N$ such that $n\geq N$ implies $\abs{a_n-L}<\eps$. Then $k\geq N$ implies $n_k\geq N$, hence $\abs{a_{n_k}-L}<\eps$.
\end{proof}

\begin{corollary}
If a sequence has two subsequences converging to distinct limits, then it is divergent.
\end{corollary}
\begin{proof}
Otherwise \cref{prop:subsequence-limit} would make both subsequences converge to the sequence's limit, contradicting \cref{thm:unique-sequence-limit}.
\end{proof}

\begin{example}
The sequence $(-1)^n$ diverges. Its even-indexed subsequence is constantly $1$, while its odd-indexed subsequence is constantly $-1$.
\end{example}

\section{Completeness Consequences}

The order completeness of the real numbers now becomes a practical tool.
It converts bounded monotone approximation into a limit, prevents Cauchy
approximations from leaving a gap, and forces bounded sequences to have
convergent subsequences.

\begin{definition}[Monotone sequence]
A sequence $(a_n)$ is \textbf{increasing} if $a_n\leq a_{n+1}$ for all $n$, \textbf{decreasing} if $a_{n+1}\leq a_n$ for all $n$, and \textbf{monotone} if it is increasing or decreasing.
\end{definition}

\begin{theorem}[Monotone convergence]
\label{thm:monotone-convergence}
An increasing sequence converges if and only if it is bounded above; in that event, its limit is $\sup\set{a_n:n\in\N}$. A decreasing sequence converges if and only if it is bounded below, and its limit is the infimum of its range.
\end{theorem}
\begin{proof}
If $a_n\to L$, choose $N_0$ so that $|a_n-L|<1$ for $n\geq N_0$.
Then $B=\max\{|L|+1,|a_1|,\ldots,|a_{N_0}|\}$ bounds every $|a_n|$. Conversely, let $s=\sup\set{a_n:n\in\N}$. Given $\eps>0$, $s-\eps$ is not an upper bound, so $a_N>s-\eps$ for some $N$. For $n\geq N$, $s-\eps<a_N\leq a_n\leq s$, and hence $\abs{a_n-s}<\eps$. The decreasing statement follows on applying the increasing statement to $(-a_n)$.
\end{proof}

The theorem immediately constructs one of the central constants of analysis.
In elementary algebra the expression \((1+1/n)^n\) appears in repeated
growth and compound-interest calculations. Here it will do more: bounded
monotone approximation will produce a real number before any exponential
function has been defined.

\begin{theorem}[The sequential construction of \(e\)]
\label{thm:e-sequence}
The sequence
\[
a_n=\left(1+\frac1n\right)^n
\qquad(n\in\N)
\]
is strictly increasing and bounded above. Consequently it converges.
\end{theorem}
\begin{proof}
\textbf{Mechanism.} Expand each finite power. For every fixed binomial
term, increasing \(n\) increases each factor, while the next expansion also
contains a new positive term.

The finite binomial identity (which follows by elementary induction) gives
\[
\begin{aligned}
a_n
&=\sum_{k=0}^{n}\binom nk\frac1{n^k}\\
&=\sum_{k=0}^{n}\frac1{k!}
  \prod_{j=0}^{k-1}\left(1-\frac jn\right),
\end{aligned}
\]
where the product for \(k=0\) is \(1\). For \(0\leq k\leq n\), every
factor in the \(k\)th term is nonnegative and is no smaller when \(n\) is
replaced by \(n+1\). The \(k=2\) term increases strictly, and the expansion
of \(a_{n+1}\) has one additional positive term. Hence \(a_{n+1}>a_n\).

Every factor in the product is at most \(1\), so
\[
a_n\leq\sum_{k=0}^{n}\frac1{k!}.
\]
For \(k\geq2\), one has \(k!\geq2\cdot3^{k-2}\). Therefore the finite
geometric sum gives
\[
a_n\leq2+\frac12\sum_{j=0}^{n-2}\frac1{3^j}<\frac{11}{4}<3.
\]
Thus \((a_n)\) is increasing and bounded above, so it converges by
\cref{thm:monotone-convergence}.
\end{proof}

\begin{definition}[The number \(e\)]
\label{def:e-sequential}
Define
\[
\boxed{
e:=\lim_{n\to\infty}\left(1+\frac1n\right)^n.
}
\]
In particular, \(2<e<3\). This definition uses only finite powers and
sequence completeness. Chapter~\ref{ch:differentiation} will construct the
exponential function independently and prove that its value at \(1\) is
this same number.
\end{definition}

\begin{example}[Recognizing the defining sequence for \(e\)]
Write $a_n=(1+1/n)^n$, so $a_n\to e$. Three altered expressions reduce
to this sequence by operations already available:
\[
\begin{aligned}
\left(1+\frac1{2n}\right)^{2n}&=a_{2n}\longrightarrow e,\\
\left(1+\frac1{n+1}\right)^n
 &=\frac{a_{n+1}}{1+\frac1{n+1}}\longrightarrow e,\\
\left(1-\frac1n\right)^{-n}
 &=a_{n-1}\left(1+\frac1{n-1}\right)\longrightarrow e\qquad(n\geq2).
\end{aligned}
\]
The first is a subsequence; the second removes one factor after an index
shift; the third rewrites a reciprocal and adds one factor. Thus recognition,
index shifts, and limit algebra do the work, with only integer powers.
Chapter~\ref{ch:differentiation} will define real powers and use logarithms
to prove the broader principle $(1+x/n)^n\to e^x$
(\cref{prop:exp-binomial-limit}). That formula is a forward pointer, not a
tool used here.
\end{example}

\begin{theorem}[Cauchy criterion]
\label{thm:cauchy-criterion}
A real sequence converges if and only if it is Cauchy.
\end{theorem}
\begin{proof}
If $a_n\to L$, choose $N$ with $\abs{a_n-L}<\eps/2$ for $n\geq N$ and, for all $m,n\geq N$, write
$|a_m-a_n|\leq|a_m-L|+|a_n-L|<\eps$. Conversely, the rational-Cauchy construction of $\R$ proves exactly this assertion in \cref{thm:real-cauchy-complete}.
\end{proof}

\begin{definition}[Peak index]
An index $n\in\N$ is a \textbf{peak index} of $(a_n)$ if
$a_n\geq a_m$ for every $m>n$.
\end{definition}
\begin{theorem}[Monotone subsequence]\label{thm:monotone-subsequence}
Every real sequence has a monotone subsequence.
\end{theorem}
\begin{proof}
\textbf{Strategy.} Either infinitely many terms are never exceeded later,
or eventually every term is exceeded by a later term.
If there are infinitely many peak indices, list them as
$n_1<n_2<\cdots$. Then $a_{n_k}\geq a_{n_{k+1}}$, so this subsequence
is nonincreasing. If there are only finitely many, choose $N$ larger
than all of them; if there are none, take $N=1$. Put $n_1=N$ and choose
$n_{k+1}$ to be the least $m>n_k$ with $a_m>a_{n_k}$. Such an index exists
because $n_k$ is not a peak. The resulting subsequence is strictly increasing.
This argument uses only order and works in any ordered field, without completeness.
\end{proof}

\begin{theorem}[Bolzano--Weierstrass]
\label{thm:bolzano-weierstrass}
Every bounded real sequence has a convergent subsequence.
\end{theorem}
\begin{proof}
By \cref{thm:monotone-subsequence}, the sequence has a monotone
subsequence. That subsequence is still bounded, so it converges by the
Monotone Convergence Theorem (\cref{thm:monotone-convergence}).
This is the useful habit of reusing a proved theorem rather than repeating
its supremum argument.
\end{proof}

A bounded sequence need not settle at one value. We can still ask for
upper and lower bounds on its eventual behavior by discarding an initial
segment and bounding the whole remaining tail. These bounds change
monotonically as the tail shrinks, so their limits exist even when the
original sequence diverges.

\begin{definition}[Limsup and liminf]
Let $(a_n)$ be bounded. For $n\in\N$, put
\[
s_n=\sup\set{a_k:k\geq n},\qquad i_n=\inf\set{a_k:k\geq n}.
\]
Define
\[
\lim_{n\to\infty}\sup a_n=\lim_{n\to\infty}s_n,
\qquad
\lim_{n\to\infty}\inf a_n=\lim_{n\to\infty}i_n.
\]
\end{definition}

\begin{theorem}[Basic limsup and liminf facts]
\label{thm:limsup-liminf}
For every bounded sequence, limsup and liminf exist and
\[
\lim_{n\to\infty}\inf a_n\leq\lim_{n\to\infty}\sup a_n.
\]
Moreover, $a_n\to L$ if and only if the two displayed limits both equal $L$.
\end{theorem}
\begin{proof}
The decreasing sequence $(s_n)$ and increasing sequence $(i_n)$ are
bounded by the same bounds as $(a_n)$, hence converge by monotone
convergence. Since $i_n\leq s_n$, order preservation orders their limits.
If $a_n\to L$, given $\eps>0$ choose $N$ with
$|a_k-L|<\eps/2$ for $k\geq N$. For every $n\geq N$, all terms in
the $n$th tail lie between $L-\eps/2$ and $L+\eps/2$. Thus
$|i_n-L|\leq\eps/2<\eps$ and $|s_n-L|\leq\eps/2<\eps$.
Conversely, if $i_n,s_n\to L$, choose $N_1,N_2$ for their errors
less than $\eps$. For $n\geq\max\{N_1,N_2\}$,
$L-\eps<i_n\leq a_n\leq s_n<L+\eps$.
\end{proof}

\paragraph{Why boundedness remains in this definition.}
\begin{example}[Tail ceilings and floors]
For $a_n=(-1)^n+1/n$, the even subsequence tends to $1$ and the odd
subsequence to $-1$. Removing finitely many terms cannot remove either
recurring behavior. The tail suprema decrease to $1$, while the tail
infima increase to $-1$. Thus $\limsup a_n=1$ and $\liminf a_n=-1$.
For $b_n=2+(-1)^n/n$, both tail barriers tend to $2$, so the sequence
converges despite oscillating. The limsup is the eventual ceiling,
not the largest early spike in the whole sequence.
\end{example}
For a bounded sequence, every subsequential limit lies between liminf and
limsup: its late terms lie between every fixed pair of tail bounds.
Each barrier itself is a subsequential limit. To reach the upper one $s$,
make two choices at stage $j$. First choose a tail index $N_j>k_{j-1}$
so late that $|s_{N_j}-s|<1/j$, where $k_0=0$. Then the definition of
supremum supplies $k_j\geq N_j$ with
$s_{N_j}-1/j<a_{k_j}\leq s_{N_j}$. Thus $k_j>k_{j-1}$ and
$|a_{k_j}-s|<2/j$, proving convergence. For the lower barrier $i$, first
choose $N_j>k_{j-1}$ with $|i_{N_j}-i|<1/j$, then choose $k_j\geq N_j$
with $i_{N_j}\leq a_{k_j}<i_{N_j}+1/j$. Again the error is less than $2/j$.
One can develop conventions for limsup and liminf of unbounded sequences,
but they are unnecessary here.  We retain real-valued suprema and infima
and use \cref{def:sequence-infinite-limit} when a sequence eventually
escapes every real bound.

\section{Equivalent Completeness Principles}

The least-upper-bound property is not the only possible formulation of
completeness.  The equivalences below explain why several apparently
different principles may all serve as the decisive feature distinguishing
the real numbers from the rationals.

\begin{lemma}[A Cauchy sequence with a convergent subsequence]
\label{lem:cauchy-convergent-subsequence}
If a Cauchy sequence has a subsequence converging to $L$, the entire
sequence converges to $L$. This holds in every ordered field.
\end{lemma}
\begin{proof}
Given $\eps>0$, choose $N$ so that $|a_m-a_n|<\eps/2$ for $m,n\geq N$.
Choose $K$ so that $|a_{n_k}-L|<\eps/2$ for $k\geq K$.
Fix $k=\max\{K,N\}$; then $n_k\geq k\geq N$.
For all $n\geq N$, the triangle inequality gives
$|a_n-L|\leq|a_n-a_{n_k}|+|a_{n_k}-L|<\eps$.
\end{proof}

\begin{theorem}[Equivalence of standard completeness principles]
\label{thm:completeness-equivalences}
Let $F$ be an Archimedean field. The following assertions are equivalent.
\begin{enumerate}
\item Every nonempty subset of $F$ having an upper bound has a least upper bound.
\item Every Cauchy sequence in $F$ converges.
\item Every bounded sequence in $F$ has a convergent subsequence.
\item Every increasing sequence in $F$ that is bounded above converges.
\item Whenever
\[
[a_1,b_1]\supseteq[a_2,b_2]\supseteq\cdots
\]
and
\[
\lim_{n\to\infty}(b_n-a_n)=0,
\]
the intervals have a common point.
\end{enumerate}
Consequently, each assertion may be used as a completeness axiom for an Archimedean field.
\end{theorem}
\begin{proof}
\textbf{Strategy.} Prove the single cycle
$(a)\Rightarrow(d)\Rightarrow(c)\Rightarrow(b)\Rightarrow(e)\Rightarrow(a)$.
The order-limit laws and the preceding lemma use only ordered-field laws,
so they apply in $F$ without assuming real completeness.

\emph{$(a)\Rightarrow(d)$.} For a bounded increasing sequence take the
supremum $s$ of its range. For $\eps>0$ some $a_N>s-\eps$, and
$s-\eps<a_N\leq a_n\leq s$ for $n\geq N$.

\emph{$(d)\Rightarrow(c)$.} The peak-index proof extracts a monotone
subsequence in $F$. It remains bounded, so (d), applied to its negative
if it is decreasing, gives convergence.

\emph{$(c)\Rightarrow(b)$.} Choose $N_0$ with $|a_m-a_n|<1$ for $m,n\geq N_0$.
For $n\geq N_0$, $|a_n|<|a_{N_0}|+1$, so
$\max\{|a_1|,\ldots,|a_{N_0}|+1\}$ bounds the whole sequence.
By (c) it has a convergent subsequence. Apply
\cref{lem:cauchy-convergent-subsequence}.

\emph{$(b)\Rightarrow(e)$.} For nested intervals with widths tending to
zero, given $\eps>0$ choose $N$ with $b_N-a_N<\eps$.
For all $m,n\geq N$, both left endpoints lie in $[a_N,b_N]$, so
$|a_m-a_n|\leq b_N-a_N<\eps$. Thus $(a_n)$ is
Cauchy; let $a_n\to x$. For each fixed $N$, the late left endpoints
belong to $[a_N,b_N]$. Order-limit preservation gives $a_N\leq x\leq b_N$.

\emph{$(e)\Rightarrow(a)$.} Let $S\ne\varnothing$ be bounded above,
choose $s_0\in S$, and choose an upper bound $u_0$. If $s_0$ is an upper
bound it is the maximum. Otherwise bisect $[s_0,u_0]$, retaining the lower
half when its midpoint is an upper bound of $S$, and the upper half otherwise.
The resulting nested intervals $[s_n,u_n]$ have $s_n$ not an upper bound,
$u_n$ an upper bound, and width $2^{-n}(u_0-s_0)$. The Archimedean property
and $2^n\geq n+1$ imply that these widths tend to zero in $F$.
By (e) choose a common point $c$. If $t\in S$ and $t>c$, choose $N_1$ with $u_{N_1}-s_{N_1}<t-c$; then
$u_{N_1}\leq c+(u_{N_1}-s_{N_1})<t$, a contradiction.
If $v<c$, choose $N_2$ with $u_{N_2}-s_{N_2}<c-v$;
then $s_{N_2}\geq c-(u_{N_2}-s_{N_2})>v$, and since
$s_{N_2}$ is not an upper bound, neither is $v$. Thus $c=\sup S$.
\end{proof}



\begin{center}
\begin{tikzpicture}[node distance=6mm]
\node (a) {least-upper-bound property};
\node (b) [below=of a] {monotone convergence};
\node (c) [below=of b] {Bolzano--Weierstrass};
\node (d) [below=of c] {Cauchy completeness};
\node (e) [below=of d] {nested intervals};
\draw[->] (a)--(b);\draw[->] (b)--(c);\draw[->] (c)--(d);\draw[->] (d)--(e);
\draw[->] (e.east)--++(1.5,0)|-(a.east);
\end{tikzpicture}
\end{center}
Each arrow has one mechanism: a supremum, peak indices, the Cauchy
estimate, shrinking widths, and finally bisection.

\section*{Exercises}
\addcontentsline{toc}{section}{Exercises}

\begin{hexercise}{seq-direct}
Prove directly that $3/(n+2)\to0$ and $(2n+1)/(n+1)\to2$. For each, show the backward discovery of $N$ separately from the forward proof.
\end{hexercise}

\begin{hexercise}{seq-laws}
Evaluate $((4k^2+1)/(2k^2-k))$, $((k+1)/(k^3+2))$, and $((-2k^4+k)/(k^2+1))$ as $k\to\infty$. Identify the degrees and leading coefficients and justify that denominators are nonzero on a tail.
\end{hexercise}

\begin{hexercise}{seq-geometric}
Give explicit thresholds proving $(3/4)^n\to0$ and $(3/2)^n\to+\infty$. Determine whether $(-3/2)^n$ has either infinite limit.
\end{hexercise}

\begin{hexercise}{seq-radical}
For $c>0$, find the limits of $\sqrt{k^2+ck}-k$ and $\sqrt{k^2+c}-k$. Explain which part of rationalization determines the different answers.
\end{hexercise}

\begin{hexercise}{seq-model}
An approximation begins with error at most $6$ units. Each correction reduces the error magnitude to at most one third of its preceding magnitude. Formulate a sequence inequality, prove convergence of the errors to zero, and find a number of corrections guaranteeing error below $1/100$ unit.
\end{hexercise}

\begin{hexercise}{seq-disprove}
Disprove the proposed limit $1$ for $a_n=2+1/n$ using the negated definition. Then prove that $2+(-1)^n$ has no limit. Explain why these are different tasks.
\end{hexercise}

\begin{exercise}
Find the limsup and liminf of $(-1)^n+2/n$ and exhibit subsequences attaining them. Prove from the tail bounds that no other subsequential limit is possible.
\end{exercise}

\begin{exercise}
Let $a_1=0$ and $a_{n+1}=(a_n+2)/3$. Prove monotonicity and boundedness, deduce convergence, and determine the limit. Which step uses completeness?
\end{exercise}

\begin{hexercise}{supremum-proof}
Let $(a_n)$ be increasing and bounded above, and let $s$ be the supremum of its range. Adapt the proof of \cref{thm:monotone-convergence} to show directly that
\[
\lim_{n\to\infty}a_n=s.
\]
Identify the exact use of the least-upper-bound property.
\end{hexercise}

\begin{exercise}
Give an example of an unbounded sequence for which
\[
\lim_{n\to\infty}\sup a_n
\]
cannot be interpreted as a real number. Explain the purpose of the postponed extended convention.
\end{exercise}

\begin{hexercise}{decreasing-completeness}
Assuming the least-upper-bound property, prove that every decreasing sequence that is bounded below converges, without appealing to the increasing case as a black box.
\end{hexercise}

\begin{exercise}
Explain why changing the diagonal digits to only $0$ and $9$ would require
an extra ambiguity check. Prove $0.12999\cdots=0.13$ from truncations.
\end{exercise}
